%%%%%%%%%%%%%%%%%%%%%%%%%%%%%%%%%%%%%%%%%
% Guide: Creating R Markdown Documents
%%%%%%%%%%%%%%%%%%%%%%%%%%%%%%%%%%%%%%%%%

\documentclass[paper=a4, fontsize=11pt]{scrartcl}

\usepackage[T1]{fontenc}
\usepackage[english]{babel}
\usepackage{amsmath,amsfonts,amssymb}
\usepackage{enumerate}
\usepackage{listings}
\usepackage[dvipsnames]{xcolor}
\usepackage{mdframed}
\usepackage{graphicx}
\usepackage{booktabs}

\usepackage{sectsty}
\allsectionsfont{\normalfont\scshape}

\usepackage[hidelinks]{hyperref}
\hypersetup{colorlinks=true, citecolor=blue, urlcolor=blue, pdfborder=000}

\usepackage{fancyhdr}
\pagestyle{fancyplain}
\fancyhead{}
\fancyfoot[L]{}
\fancyfoot[C]{}
\fancyfoot[R]{\thepage}
\renewcommand{\headrulewidth}{0pt}
\renewcommand{\footrulewidth}{0pt}
\setlength{\headheight}{13.6pt}

\setlength\parindent{0pt}

% Horizontal rule
\newcommand{\horrule}[1]{\rule{\linewidth}{#1}}

% R code styling
\lstdefinestyle{rcode}{
  backgroundcolor=\color{gray!8},
  basicstyle=\ttfamily\small,
  breaklines=true,
  frame=single,
  rulecolor=\color{gray!40},
  framesep=4pt,
  xleftmargin=6pt,
  xrightmargin=6pt,
  language=R,
  keywordstyle=\color{NavyBlue},
  stringstyle=\color{Maroon},
  commentstyle=\color{OliveGreen},
}

\lstdefinestyle{yaml}{
  backgroundcolor=\color{gray!8},
  basicstyle=\ttfamily\small,
  breaklines=true,
  frame=single,
  rulecolor=\color{gray!40},
  framesep=4pt,
  xleftmargin=6pt,
  xrightmargin=6pt,
}

\lstdefinestyle{shell}{
  backgroundcolor=\color{black!90},
  basicstyle=\ttfamily\small\color{white},
  breaklines=true,
  frame=single,
  framesep=4pt,
  xleftmargin=6pt,
  xrightmargin=6pt,
}

% Tip / warning box
\newmdenv[
  backgroundcolor=blue!5,
  linecolor=blue!40,
  linewidth=1pt,
  leftmargin=0pt, rightmargin=0pt,
  innerleftmargin=8pt, innerrightmargin=8pt,
  innertopmargin=6pt, innerbottommargin=6pt,
]{tipbox}

\newmdenv[
  backgroundcolor=orange!8,
  linecolor=orange!60,
  linewidth=1pt,
  leftmargin=0pt, rightmargin=0pt,
  innerleftmargin=8pt, innerrightmargin=8pt,
  innertopmargin=6pt, innerbottommargin=6pt,
]{warnbox}

%----------------------------------------------------------------------------------------

\title{
\normalfont \normalsize
\horrule{0.5pt} \\[0.4cm]
\huge A Step-by-Step Guide to \\[0.2cm] R Markdown (\texttt{.Rmd}) \\
\horrule{2pt} \\[0.5cm]
}

\author{Prof.\ Tzu-Ting Yang}
\date{\normalsize\today}

\begin{document}
\maketitle

\tableofcontents
\bigskip\horrule{0.4pt}\bigskip

%----------------------------------------------------------------------------------------
\section{What Is R Markdown?}
%----------------------------------------------------------------------------------------

An \textbf{R Markdown} file (\texttt{.Rmd}) is a plain-text document that combines
\textit{narrative text} (written in Markdown) with \textit{R code chunks}.
When you \textit{knit} it, R runs all the code, captures every output and figure,
and the \texttt{rmarkdown} package converts the result into a polished
\textbf{HTML} (or PDF, Word) report.

\medskip
The workflow is:
\[
\texttt{.Rmd file}
\;\xrightarrow{\texttt{knitr}}\;
\text{R runs code} + \text{captures output}
\;\xrightarrow{\texttt{rmarkdown}}\;
\texttt{.html}
\]

\begin{tipbox}
\textbf{Advantage over Stata Markdown:} R Markdown does \textit{not} require a
separate Pandoc installation --- the \texttt{rmarkdown} package bundles Pandoc
automatically. Setup is simpler: install R, RStudio, and two packages.
\end{tipbox}

\bigskip\horrule{0.4pt}\bigskip

%----------------------------------------------------------------------------------------
\section{Course Templates on Dropbox}
%----------------------------------------------------------------------------------------

\textbf{Goal:} produce a single \texttt{.html} file that records all your
code, output, and figures in one document.

\medskip
Complete working templates are available in the \textbf{course Dropbox folder}.
Two sub-folders are relevant:

\medskip
\begin{center}
\begin{tabular}{p{4.5cm} p{9.5cm}}
\toprule
\textbf{Folder} & \textbf{Contents} \\
\midrule
\texttt{R\_code/} &
Regular \texttt{.R} scripts for each method.
Use these to understand the R code before converting to \texttt{.Rmd}. \\[4pt]
\texttt{R\_code/R\_Markdown/} &
Ready-to-knit \texttt{.Rmd} templates --- one per method.
These are your \textbf{main reference} when writing your own \texttt{.Rmd} file. \\
\bottomrule
\end{tabular}
\end{center}

\medskip
The available templates in \texttt{R\_code/R\_Markdown/} are:

\medskip
\begin{center}
\begin{tabular}{ll}
\toprule
\textbf{File} & \textbf{Method} \\
\midrule
\texttt{regression.Rmd}     & OLS regression and controls \\
\texttt{DID.Rmd}            & Difference-in-Differences \\
\texttt{Stag\_DID.Rmd}      & Staggered DID \\
\texttt{Syn\_DID.Rmd}       & Synthetic DID \\
\texttt{fixed\_effects.Rmd} & Fixed effects \\
\texttt{matching.Rmd}       & Matching estimators \\
\texttt{RDD.Rmd}            & Regression Discontinuity Design \\
\texttt{IV.Rmd}             & Instrumental Variables \\
\texttt{SSIV.Rmd}           & Shift-Share IV \\
\texttt{SCM.Rmd}            & Synthetic Control Method \\
\texttt{ML.Rmd}             & Machine Learning \\
\bottomrule
\end{tabular}
\end{center}

\begin{tipbox}
\textbf{Recommended workflow:}
\begin{enumerate}
  \item Open the \texttt{.Rmd} template closest to your method (e.g.\ \texttt{DID.Rmd}).
  \item Read how the YAML header, path setup chunk, and \texttt{\{r\}} code chunks are structured.
  \item Create your own \texttt{.Rmd} file following the same structure, replacing the
        example data and code with your own.
  \item Click \textbf{Knit} in RStudio (or run \texttt{rmarkdown::render()}) to produce
        \texttt{yourfile.html}.
\end{enumerate}
\end{tipbox}

\bigskip\horrule{0.4pt}\bigskip

%----------------------------------------------------------------------------------------
\section{What You Need to Install (One-Time Setup)}
%----------------------------------------------------------------------------------------

\begin{center}
\begin{tabular}{clll}
\toprule
\textbf{Step} & \textbf{What} & \textbf{Where} & \textbf{Required?} \\
\midrule
1 & R (base language) & \url{https://cran.r-project.org} & Yes \\
2 & RStudio (IDE)     & \url{https://posit.co/download/rstudio-desktop} & Strongly recommended \\
3 & \texttt{rmarkdown} R package & Inside R & Yes \\
4 & \texttt{knitr} R package     & Inside R & Yes \\
\bottomrule
\end{tabular}
\end{center}

\begin{tipbox}
\textbf{No separate Pandoc install needed!} Unlike Stata Markdown,
\texttt{rmarkdown} bundles its own Pandoc. Once you install \texttt{rmarkdown},
you are ready to compile HTML reports immediately.
\end{tipbox}

\bigskip\horrule{0.4pt}\bigskip

%----------------------------------------------------------------------------------------
\section{Step 1 \& 2 --- Install R and RStudio}
%----------------------------------------------------------------------------------------

\begin{enumerate}
  \item \textbf{Install R}: go to \url{https://cran.r-project.org}, choose your OS,
        download and run the installer. Accept all defaults.
  \item \textbf{Install RStudio}: go to \url{https://posit.co/download/rstudio-desktop},
        download the free Desktop version, and install it.
\end{enumerate}

\begin{tipbox}
Always install \textbf{R first}, then RStudio. RStudio detects R automatically during
installation.
\end{tipbox}

\bigskip\horrule{0.4pt}\bigskip

%----------------------------------------------------------------------------------------
\section{Step 3 \& 4 --- Install R Packages}
%----------------------------------------------------------------------------------------

Open RStudio and run the following once in the \textbf{Console}:

\begin{lstlisting}[style=rcode]
install.packages(c("rmarkdown", "knitr"))
\end{lstlisting}

Also install the analysis packages used throughout the course:

\begin{lstlisting}[style=rcode]
install.packages(c(
  "haven",     # read Stata .dta files
  "dplyr",     # data manipulation
  "ggplot2",   # graphs
  "fixest",    # fast fixed effects / IV regression
  "estimatr",  # robust standard errors
  "broom",     # tidy regression output
  "modelsummary" # publication-quality regression tables
))
\end{lstlisting}

Verify the installation:

\begin{lstlisting}[style=rcode]
library(rmarkdown)   # should load silently
rmarkdown::pandoc_version()  # prints Pandoc version bundled with rmarkdown
\end{lstlisting}

\bigskip\horrule{0.4pt}\bigskip

%----------------------------------------------------------------------------------------
\section{Step 5 --- Write Your \texttt{.Rmd} File}
%----------------------------------------------------------------------------------------

An \texttt{.Rmd} file has three parts: a YAML header, Markdown narrative text,
and R code chunks.

\subsection*{5a. YAML Header}

The YAML header sits between two \texttt{-{}-{}-} lines at the very top of the file.
It controls the document title, author, date, and output format.

\begin{lstlisting}[style=yaml]
---
title: "My Empirical Report"
author: "Your Name"
date: "2025"
output:
  html_document:
    toc: true          # table of contents
    toc_float: true    # floating TOC sidebar
    theme: flatly      # Bootstrap theme
    highlight: tango   # code syntax highlighting
    number_sections: true
---
\end{lstlisting}

\subsection*{5b. Setup Chunk}

Place this immediately after the YAML header.
It sets global chunk options and configures file paths.

\begin{lstlisting}[style=rcode]
```{r setup, include=FALSE}
knitr::opts_chunk$set(echo = TRUE, warning = FALSE, message = FALSE)

rm(list = ls())

username <- Sys.info()[["user"]]
if (username == "alice") {
  rawdata <- "C:/Users/alice/Desktop/data"
} else if (username == "bob") {
  rawdata <- "/Users/bob/Documents/data"
}
```
\end{lstlisting}

\texttt{include=FALSE} means the chunk runs but neither the code nor its output
appears in the HTML. Find your username in R with \texttt{Sys.info()[["user"]]}.

\subsection*{5c. Install Packages (show but do not run)}

\begin{lstlisting}[style=rcode]
```{r packages, eval=FALSE}
install.packages(c("haven", "ggplot2", "dplyr", "fixest"))
```
\end{lstlisting}

\texttt{eval=FALSE} displays the code in the report but does \textit{not} run it
during knitting --- important, because \texttt{install.packages()} should never run
automatically.

\subsection*{5d. Load Libraries}

\begin{lstlisting}[style=rcode]
```{r libraries}
library(haven)
library(ggplot2)
library(dplyr)
library(fixest)
```
\end{lstlisting}

\subsection*{5e. Analysis Sections}

Use Markdown headings (\texttt{\#}, \texttt{\#\#}) between chunks to add narrative:

\begin{lstlisting}[style=yaml]
# 1. Load Data

```{r load}
df <- read_dta(file.path(rawdata, "mydata.dta"))
summary(df)
```

# 2. Regression

```{r reg}
model <- feols(y ~ treat + x1 + x2 | id + year, data = df)
summary(model)
```
\end{lstlisting}

\bigskip\horrule{0.4pt}\bigskip

%----------------------------------------------------------------------------------------
\section{Chunk Options Reference}
%----------------------------------------------------------------------------------------

\begin{center}
\begin{tabular}{lll p{6cm}}
\toprule
\textbf{Option} & \textbf{Values} & \textbf{Default} & \textbf{Effect} \\
\midrule
\texttt{echo}    & TRUE/FALSE & TRUE  & Show the R code in the report \\
\texttt{eval}    & TRUE/FALSE & TRUE  & Run the code \\
\texttt{include} & TRUE/FALSE & TRUE  & Include any output at all \\
\texttt{warning} & TRUE/FALSE & TRUE  & Show R warning messages \\
\texttt{message} & TRUE/FALSE & TRUE  & Show R messages (e.g.\ from \texttt{library()}) \\
\texttt{fig.width} & number & 7 & Figure width in inches \\
\texttt{fig.height} & number & 5 & Figure height in inches \\
\bottomrule
\end{tabular}
\end{center}

\medskip
Recommended global defaults (in the setup chunk):

\begin{lstlisting}[style=rcode]
knitr::opts_chunk$set(echo = TRUE, warning = FALSE, message = FALSE)
\end{lstlisting}

\bigskip\horrule{0.4pt}\bigskip

%----------------------------------------------------------------------------------------
\section{Step 6 --- Compile to HTML}
%----------------------------------------------------------------------------------------

\subsection*{Method A --- Knit Button (RStudio)}

Click the \textbf{Knit} button at the top of the editor (blue yarn icon).
RStudio saves the \texttt{.html} file in the same folder as your \texttt{.Rmd} and
opens it automatically in the Viewer.

\subsection*{Method B --- R Console}

\begin{lstlisting}[style=rcode]
rmarkdown::render("myreport.Rmd")
\end{lstlisting}

Or with an explicit output path:

\begin{lstlisting}[style=rcode]
rmarkdown::render("myreport.Rmd",
                  output_file = "myreport.html",
                  output_dir  = "C:/Users/YourName/output")
\end{lstlisting}

\begin{tipbox}
The HTML output is \textbf{self-contained by default} --- all figures are embedded
inline. You do not need a separate \texttt{bundle} option like Stata Markdown.
Simply submit the single \texttt{.html} file.
\end{tipbox}

\bigskip\horrule{0.4pt}\bigskip

%----------------------------------------------------------------------------------------
\section{Troubleshooting Common Errors}
%----------------------------------------------------------------------------------------

\begin{center}
\begin{tabular}{p{5cm} p{9cm}}
\toprule
\textbf{Error / Problem} & \textbf{Cause and Fix} \\
\midrule
\texttt{could not find function "feols"} &
Package not loaded. Add \texttt{library(fixest)} at the top of your \texttt{.Rmd}. \\[6pt]

\texttt{there is no package called 'X'} &
Package not installed. Run \texttt{install.packages("X")} in the Console (not in a chunk with \texttt{eval=TRUE}). \\[6pt]

\texttt{cannot open file: No such file} &
Path is wrong or \texttt{rawdata} not set. Check your path setup chunk and \texttt{Sys.info()[["user"]]}. \\[6pt]

Figures not showing in HTML &
Make sure the chunk that generates the plot does not have \texttt{include=FALSE}. \\[6pt]

Knit fails at \texttt{install.packages()} &
Always use \texttt{eval=FALSE} on install chunks. \texttt{install.packages()} must not run during knitting. \\[6pt]

Code output missing &
Check if \texttt{echo=FALSE} is set globally or on that chunk.
Remove it or set \texttt{echo=TRUE}. \\

\bottomrule
\end{tabular}
\end{center}

\bigskip\horrule{0.4pt}\bigskip

%----------------------------------------------------------------------------------------
\section{Quick Reference: Complete Setup Checklist}
%----------------------------------------------------------------------------------------

\begin{enumerate}
  \item[\(\square\)] Install \textbf{R} from \url{https://cran.r-project.org}
  \item[\(\square\)] Install \textbf{RStudio} from \url{https://posit.co/download/rstudio-desktop}
  \item[\(\square\)] In RStudio Console: \texttt{install.packages(c("rmarkdown","knitr"))}
  \item[\(\square\)] Install analysis packages: \texttt{haven}, \texttt{dplyr}, \texttt{ggplot2}, \texttt{fixest}, etc.
  \item[\(\square\)] Verify: \texttt{rmarkdown::pandoc\_version()} prints a version number
  \item[\(\square\)] Write your \texttt{.Rmd} file: YAML header + setup chunk + \texttt{\{r\}} chunks
  \item[\(\square\)] Set paths using \texttt{Sys.info()[["user"]]} in the setup chunk
  \item[\(\square\)] Click \textbf{Knit} (or run \texttt{rmarkdown::render("file.Rmd")})
  \item[\(\square\)] Open \texttt{file.html} in any browser --- done!
\end{enumerate}

\bigskip\horrule{0.4pt}\bigskip

%----------------------------------------------------------------------------------------
\section{A Complete \texttt{.Rmd} Template}
%----------------------------------------------------------------------------------------

Copy and save the following as \texttt{myreport.Rmd}.

\begin{lstlisting}[style=yaml]
---
title: "Your Report Title"
author: "Your Name"
date: "2025"
output:
  html_document:
    toc: true
    toc_float: true
    theme: flatly
    highlight: tango
    number_sections: true
---

```{r setup, include=FALSE}
knitr::opts_chunk$set(echo = TRUE, warning = FALSE, message = FALSE)
rm(list = ls())

username <- Sys.info()[["user"]]
if (username == "alice") {
  rawdata <- "C:/Users/alice/Desktop/data"
} else if (username == "bob") {
  rawdata <- "/Users/bob/Documents/data"
}
```

# 1. Install Packages (run once)

```{r packages, eval=FALSE}
install.packages(c("haven", "dplyr", "ggplot2", "fixest"))
```

# 2. Load Libraries

```{r libraries}
library(haven)
library(dplyr)
library(ggplot2)
library(fixest)
```

# 3. Load Data

```{r load}
df <- read_dta(file.path(rawdata, "mydata.dta"))
summary(df)
```

# 4. Visualize

```{r plot, fig.width=7, fig.height=4}
ggplot(df, aes(x = treatment, y = outcome)) +
  geom_boxplot() +
  theme_minimal()
```

# 5. Regression

```{r reg}
model <- feols(outcome ~ treatment + control1 + control2 | id + year, data = df)
summary(model)
```
\end{lstlisting}

\medskip
Knit with the button in RStudio, or from the Console:

\begin{lstlisting}[style=rcode]
rmarkdown::render("myreport.Rmd")
\end{lstlisting}

\end{document}
